% ================================================================
% THE HEISENBERG ROSETTA STONE
% ================================================================
% This chapter lives in BOTH volumes.
% Volume I: chapters/theory/rosetta_stone.tex
% Volume II: chapters/examples/rosetta_stone.tex
%
% It proves, by explicit computation on H_k, that the bar complex
% is a Swiss-cheese algebra in the sense of Definition~\ref{def:SC}.
% The chapter uses only macros available in both volumes.
% ================================================================

% Fallback macros: ensure compilation in both volumes.
\providecommand{\barB}{\bar{B}}
\providecommand{\cA}{\mathcal{A}}
\providecommand{\cH}{\mathcal{H}}
\providecommand{\SCchtop}{\mathsf{SC}^{\mathrm{ch,top}}}
\providecommand{\Ainf}{\mathsf{A}_{\infty}}
\providecommand{\Eone}{\mathsf{E}_1}
\providecommand{\Einf}{\mathsf{E}_{\infty}}
\providecommand{\Ass}{\mathsf{Ass}}
\providecommand{\FM}{\mathrm{FM}}
\providecommand{\Conf}{C}
\providecommand{\dfib}{d_{\mathrm{fib}}}
\providecommand{\Dg}[1]{D_{#1}}
\providecommand{\dzero}{d_0}
\providecommand{\barBch}{\bar{B}^{\mathrm{ch}}}
\providecommand{\Omegach}{\Omega^{\mathrm{ch}}}
\providecommand{\bC}{\mathbb{C}}
\providecommand{\bR}{\mathbb{R}}
\providecommand{\C}{\mathbb{C}}
\providecommand{\R}{\mathbb{R}}
\providecommand{\Z}{\mathbb{Z}}
\providecommand{\fg}{\mathfrak{g}}
\providecommand{\id}{\mathrm{id}}
\providecommand{\Res}{\operatorname{Res}}
\providecommand{\Ran}{\operatorname{Ran}}
\providecommand{\chirCom}{\mathsf{Com}^{\mathrm{ch}}}
\providecommand{\chirLie}{\mathsf{Lie}^{\mathrm{ch}}}
\providecommand{\Linf}{\mathsf{L}_{\infty}}
\providecommand{\ClaimStatusProvedHere}{\textnormal{[proved here]}}
\providecommand{\ClaimStatusProvedElsewhere}{\textnormal{[proved elsewhere]}}
\providecommand{\ClaimStatusConjectured}{\textnormal{[conjectured]}}
\providecommand{\ClaimStatusHeuristic}{\textnormal{[physical heuristic]}}
\providecommand{\Hom}{\operatorname{Hom}}
\providecommand{\colldiv}[2]{D_{#1#2}}
\providecommand{\LogForm}[2]{\eta_{#1#2}}


\section{The Heisenberg Rosetta Stone}
\label{sec:rosetta-stone}
\index{Rosetta Stone|textbf}
\index{Heisenberg algebra!Rosetta Stone|textbf}
\index{Swiss-cheese operad!Heisenberg example|textbf}

As the Steinberg variety presents the Hecke algebra through a single
correspondence---flag variety on one side, nilpotent cone on the
other, convolution algebra in between---the Heisenberg bar complex
presents the $\Ainf$ chiral algebra through a single dg coalgebra:
the bar differential sees the holomorphic direction, the coproduct
sees the topological direction, and the combined Swiss-cheese structure
is the entire content of the 3d theory on $\bC \times \bR$.

The present section proves this for the Heisenberg algebra
$\cH_k$ by explicit computation.  The results are neither deep
nor difficult; they are \emph{explicit}, and the explicitness
carries information that no amount of abstract machinery can supply.
The Heisenberg algebra has a single generator~$J$, a single
OPE pole of order~$2$, and no composite fields.  These three
simplifications---one generator, one pole order, no composites---make
every step checkable by hand.  Yet none of the five main theorems
of Volume~I becomes trivial: bar-cobar adjunction (Theorem~A),
inversion (Theorem~B), complementarity (Theorem~C), the modular
characteristic (Theorem~D), and Hochschild cohomology (Theorem~H)
are all instantiated here with full content.

The section is organized around five revelations.  The first
(\S\ref{subsec:rosetta-swiss-cheese}) identifies the bar complex
as a Swiss-cheese coalgebra; the second
(\S\ref{subsec:rosetta-operations}) reads off all higher operations
and observes that they vanish; the third
(\S\ref{subsec:rosetta-pva}) extracts the Poisson vertex algebra
on cohomology; the fourth (\S\ref{subsec:rosetta-genus1}) watches
curvature appear at genus~$1$; and the fifth
(\S\ref{subsec:rosetta-bridge}) assembles a dictionary that
maps every bar-complex datum to both a Volume~I theorem and a
Volume~II physical identification.


% ================================================================
\subsection{The bar complex as a Swiss-cheese algebra}
\label{subsec:rosetta-swiss-cheese}
\index{bar complex!Swiss-cheese structure!Heisenberg}
% ================================================================

We fix notation.  Let $V = \bC \cdot J$ be the one-dimensional
vector space spanned by the Heisenberg current, let $k \in \bC$
be the level, and let $\cH_k$ be the corresponding Heisenberg
chiral algebra with OPE
\begin{equation}\label{eq:rosetta-heisenberg-ope}
J(z)\, J(w) \;\sim\; \frac{k}{(z - w)^2}.
\end{equation}
The bar complex $\barB^{\mathrm{ch}}(\cH_k)$ is the cofree
conilpotent coalgebra on the desuspension $s^{-1}\overline{\cH}_k$,
equipped with the bar differential $d_{\barB}$ constructed from
OPE residues on the Fulton--MacPherson compactification.
Concretely, a bar element of tensor degree~$n$ lives on
\begin{equation}\label{eq:rosetta-bar-element-space}
\barB^n(\cH_k) \;=\;
(s^{-1}\overline{\cH}_k)^{\otimes n}
\;\otimes\;
\Omega^{\bullet}\bigl(\overline{\Conf}_{n+1}(X)\bigr),
\end{equation}
where $\overline{\Conf}_{n+1}(X) = \FM_{n+1}(X)$ is the
Fulton--MacPherson compactification and the logarithmic forms
$\eta_{ij} = d\log(z_i - z_j)$ provide the propagators.

\begin{theorem}[Heisenberg bar complex is a Swiss-cheese algebra;
\ClaimStatusProvedHere]
\label{thm:rosetta-swiss-cheese}
\index{Swiss-cheese operad!Heisenberg bar complex}
The geometric bar complex $\barB^{\mathrm{ch}}(\cH_k)$ with:
\begin{enumerate}[label=\textup{(\roman*)}]
\item bar differential $d_{\barB}$
  \textup{(}residues on collision divisors of\/ $\FM_n(\bC)$\textup{)}
  as closed-color structure,
\item deconcatenation coproduct $\Delta$
  \textup{(}ordered splitting\textup{)} as open-color structure,
\end{enumerate}
is an algebra over $\SCchtop$.
\end{theorem}

\begin{proof}
The proof has four steps, each corresponding to one axiom of the
Swiss-cheese operad.

\medskip
\noindent\textit{Step~1: Closed color.}\enspace
The bar differential $d_{\barB}$ is the coderivation of the
cofree coalgebra $T^c(s^{-1}\overline{\cH}_k)$ determined by
the OPE residue.  For the Heisenberg algebra, the only nonvanishing
component is the binary one:
\begin{equation}\label{eq:rosetta-dbar-binary}
d_{\barB}\bigl[s^{-1}J \,|\, s^{-1}J\bigr]
\;=\;
\Res_{z_1 = z_2}\!\left[
  J(z_1)\, J(z_2) \cdot \eta_{12}
\right]
\;=\;
\Res_{z_1 = z_2}\!\left[
  \frac{k}{(z_1 - z_2)^2} \cdot d\log(z_1 - z_2)
\right].
\end{equation}
The residue extracts the coefficient of $(z_1 - z_2)^{-1}$
in the integrand $k(z_1 - z_2)^{-2} \cdot (z_1 - z_2)^{-1}\, dz_1
= k(z_1 - z_2)^{-3}\, dz_1$, which gives
$d_{\barB}[s^{-1}J \,|\, s^{-1}J] = k$.

The nilpotence $d_{\barB}^2 = 0$ is the Arnold relation
on $\FM_3(\bC)$.  The three codimension-$1$ boundary strata of
$\FM_3(\bC)$ are the collision divisors $\colldiv{1}{2}$,
$\colldiv{2}{3}$, and $\colldiv{1}{3}$.
On these strata, the logarithmic $2$-form
$\eta_{12} \wedge \eta_{23}$ restricts to the Arnold relation
\begin{equation}\label{eq:rosetta-arnold}
[\eta_{12} \wedge \eta_{23}]
\;+\;
[\eta_{23} \wedge \eta_{13}]
\;+\;
[\eta_{13} \wedge \eta_{12}]
\;=\; 0
\end{equation}
in $H^2(\FM_3(\bC))$.
Each summand contributes a residue to $d_{\barB}^2$, and their
sum vanishes by exactness.  For the Heisenberg algebra, this
computation produces $d_{\barB}^2[s^{-1}J \,|\, s^{-1}J \,|\, s^{-1}J]
= k \cdot (\text{Arnold sum}) = 0$.

The closed-color structure is thus a dg coalgebra
$(\barB^{\mathrm{ch}}(\cH_k), d_{\barB})$, with $d_{\barB}$
built from $\FM_n(\bC)$.  This is holomorphic factorization.

\medskip
\noindent\textit{Step~2: Open color.}\enspace
The deconcatenation coproduct
\begin{equation}\label{eq:rosetta-delta}
\Delta\bigl[a_1 \,|\, \cdots \,|\, a_n\bigr]
\;=\;
\sum_{i=0}^{n}
\bigl[a_1 \,|\, \cdots \,|\, a_i\bigr]
\;\otimes\;
\bigl[a_{i+1} \,|\, \cdots \,|\, a_n\bigr]
\end{equation}
is coassociative by construction: it is the tensor coalgebra
coproduct.  The ordering of tensor factors corresponds to the
ordering of points $t_1 < t_2 < \cdots < t_n$ on $\bR$;
splitting at position~$i$ is cutting the interval
$[t_1, t_n] \subset \bR$ at the gap between $t_i$ and $t_{i+1}$.
The ordered configuration space $\Conf_n^{<}(\bR)$ is
contractible (it is an open simplex
$\{t_1 < \cdots < t_n\} \cong \bR^{n-1}_{> 0}$), so the
$\Eone$-structure on $\Conf_n(\bR)$ contracts to
$\Ass$-structure: $\Delta$ is strictly coassociative,
not merely $\Eone$-coassociative up to homotopy.

\medskip
\noindent\textit{Step~3: Compatibility.}\enspace
The coderivation property
\begin{equation}\label{eq:rosetta-compatibility}
\Delta \circ d_{\barB}
\;=\;
(\id \otimes d_{\barB} \;+\; d_{\barB} \otimes \id)
\circ \Delta
\end{equation}
asserts that $d_{\barB}$ is a coderivation of $\Delta$.
We verify this on bar elements of degree~$2$.  The left-hand side
gives
\[
\Delta\bigl(d_{\barB}[s^{-1}J \,|\, s^{-1}J]\bigr)
\;=\;
\Delta(k)
\;=\;
k \otimes [] + [] \otimes k,
\]
where $[]$ denotes the empty bar element (the coaugmentation).
The right-hand side gives
\[
(d_{\barB} \otimes \id + \id \otimes d_{\barB})\bigl(
  [] \otimes [s^{-1}J \,|\, s^{-1}J]
  \;+\;
  [s^{-1}J] \otimes [s^{-1}J]
  \;+\;
  [s^{-1}J \,|\, s^{-1}J] \otimes []
\bigr).
\]
The middle term $[s^{-1}J] \otimes [s^{-1}J]$ contributes nothing
because $d_{\barB}$ on a single element $[s^{-1}J]$ requires a
binary collision, which is absent.  The first and last terms
contribute $[] \otimes k + k \otimes []$, matching the left-hand side.

The structural reason for this agreement is geometric:
the collision divisors $\colldiv{i}{j} \subset \FM_n(\bC)$
parametrize collisions in the holomorphic direction, while the
splitting $\Delta$ acts in the topological direction on $\bR$.
These directions are independent in the product
$\FM_n(\bC) \times \Conf_n^{<}(\bR)$: a collision
$z_i \to z_j$ in $\bC$ does not merge the corresponding
points $t_i, t_j \in \bR$.  The residue operation (extracting the
singular part of the OPE along $\colldiv{i}{j}$) therefore
commutes with the splitting operation (cutting the interval at a
gap in the $\bR$-ordering).

\medskip
\noindent\textit{Step~4: Coherences.}\enspace
The Boardman--Vogt $W$-construction
$W(\SCchtop)$ replaces strict compatibility with a family of
homotopy-coherent conditions.  For $(\barB^{\mathrm{ch}}(\cH_k),
d_{\barB}, \Delta)$, all coherence conditions are automatically
satisfied at the strict level:
\begin{enumerate}[label=(\alph*)]
\item the closed-color differential $d_{\barB}$ is strictly
  square-zero (Arnold relation, not merely $\Ainf$-null);
\item the open-color coproduct $\Delta$ is strictly coassociative
  (tensor coalgebra, not merely $\Eone$-coassociative);
\item compatibility~\eqref{eq:rosetta-compatibility} holds strictly,
  not up to a homotopy parametrized by the $W$-construction.
\end{enumerate}
The $W$-construction coherences between $\FM_n(\bC)$ and
$\Conf_n^{<}(\bR)$ are trivial because the product
$\FM_n(\bC) \times \Conf_n^{<}(\bR)$ decomposes: the FM operad is
formal at genus~$0$ (Kontsevich), so all higher coherences on the
holomorphic factor vanish; the ordered configuration space on $\bR$
is contractible, so all higher coherences on the topological factor
vanish; and the product structure admits no mixed higher coherences.
\end{proof}


% ================================================================
\subsection{Reading off the higher operations}
\label{subsec:rosetta-operations}
\index{higher operations!from bar complex}
\index{Heisenberg algebra!higher operations}
% ================================================================

The Swiss-cheese structure of the bar complex carries all the
information of an $\Ainf$ chiral algebra.  The following
proposition makes the extraction explicit.

\begin{proposition}[Higher operations from bar elements;
\ClaimStatusProvedHere]
\label{prop:rosetta-operations}
\index{A-infinity operations@$\Ainf$ operations!from bar complex}
The $\Ainf$ operations $m_k$ of a chiral algebra $\cA$ are
recovered from the bar complex by the formula
\begin{equation}\label{eq:rosetta-mk-extraction}
m_k(a_1, \ldots, a_k;\, \lambda_1, \ldots, \lambda_{k-1})
\;=\;
\text{coefficient of }
s^{-1}a_1 \otimes \cdots \otimes s^{-1}a_k
\text{ in }
d_{\barB}^{(k)},
\end{equation}
where $d_{\barB}^{(k)}$ denotes the component of the bar
differential that maps $\barB^k \to \barB^{k-1}$ and the
spectral parameters $\lambda_j$ arise from the regularization
of the collision integral on $\FM_k(\bC) \times \Conf_k^{<}(\bR)$.
\end{proposition}

\begin{proof}
This is the defining property of the bar construction as a
coderivation: $d_{\barB}$ is the unique coderivation of
$T^c(s^{-1}\overline{\cA})$ whose composition with the
projection $T^c(s^{-1}\overline{\cA}) \twoheadrightarrow
s^{-1}\overline{\cA}$ recovers the sequence of operations
$(m_1, m_2, m_3, \ldots)$.  The spectral parameter
$\lambda_j$ arises from the regularized integral
$\int_{\Conf_k^{<}(\bR)} dt_1 \cdots dt_{k-1}$
along the topological direction, with $\lambda_j = t_{j+1} - t_j$
recording the separation between consecutive points.
\end{proof}

We now compute all $m_k$ for $\cH_k$.

\begin{proposition}[Heisenberg higher operations;
\ClaimStatusProvedHere]
\label{prop:rosetta-heisenberg-mk}
\index{Heisenberg algebra!$m_k$ operations}
For the Heisenberg algebra $\cH_k$:
\begin{enumerate}[label=\textup{(\roman*)}]
\item $m_1 = 0$ \textup{(}no internal differential\textup{)}.
\item $m_2(J, J;\, \lambda) = k\lambda$
  \textup{(}the OPE residue\textup{)}.
\item $m_k = 0$ for all $k \geq 3$
  \textup{(}no ternary or higher operations\textup{)}.
\end{enumerate}
In particular, $\cH_k$ is \emph{formal}: its $\Ainf$ structure is
concentrated in $m_2$.
\end{proposition}

\begin{proof}
(i) The internal differential of $\cH_k$ is zero: there is no
$Q$-operator, and the algebra has trivial cohomological grading
concentrated in degree~$0$.

\smallskip
(ii) The binary operation is the double-pole residue.  From
equation~\eqref{eq:rosetta-dbar-binary}, the bar differential
on a degree-$2$ element is
\[
d_{\barB}[s^{-1}J \,|\, s^{-1}J]
\;=\;
\Res_{z_1 = z_2}\!\left[
  \frac{k}{(z_1 - z_2)^2}
  \cdot d\log(z_1 - z_2)
\right]
\;=\; k.
\]
In the spectral (Fourier) parametrization, the meromorphic function
$k/(z_1 - z_2)^2$ corresponds to the spectral expression
$m_2(J, J;\, \lambda) = k\lambda$.  The spectral parameter
$\lambda$ records the weight of the residue: a pole of order~$2$
at $z_1 = z_2$ produces a first-order polynomial in~$\lambda$
via the substitution $(z_1 - z_2)^{-n-1} \leftrightarrow
\lambda^{(n)}/n!$.

\smallskip
(iii) The vanishing of $m_k$ for $k \geq 3$ is the decisive
computation.  Consider the degree-$3$ bar element
\[
\xi_3
\;=\;
s^{-1}J \otimes s^{-1}J \otimes s^{-1}J
\;\otimes\;
\eta_{12} \wedge \eta_{23}
\;\in\;
\barB^3(\cH_k).
\]
To compute $d_{\barB}^{(3)}$, one must extract a residue
from the iterated OPE of three Heisenberg currents along a
codimension-$2$ stratum of $\FM_3(\bC)$.  The iterated OPE is
\begin{align}\label{eq:rosetta-iterated-ope}
J(z_1)\, J(z_2)\, J(z_3)
&\;\sim\;
\frac{k}{(z_1 - z_2)^2}\, J(z_3)
\;+\;
\frac{k}{(z_1 - z_3)^2}\, J(z_2)
\;+\;
\frac{k}{(z_2 - z_3)^2}\, J(z_1).
\end{align}
Each term has a double pole in one pair of variables and is
regular in the third.  The ternary operation $m_3$ would arise
from a \emph{simultaneous} triple collision $z_1 \to z_2 \to z_3$.
On $\FM_3(\bC)$, the triple collision stratum is a point
(the maximal boundary stratum).  The residue at this point requires
a pole of order $\geq 3$ in the combined variables.  But the
iterated OPE~\eqref{eq:rosetta-iterated-ope} has at most
double poles in each pair --- there is no pole of the required
order in the ternary collision limit.

More precisely: in the coordinates $u = z_1 - z_2$,
$v = z_2 - z_3$ on $\FM_3(\bC)$, the triple collision
$u, v \to 0$ has codimension~$2$.  The form
$\eta_{12} \wedge \eta_{23} = d\log u \wedge d\log v$
contributes a factor $u^{-1}v^{-1}\, du \wedge dv$.  The
OPE contributes at most $u^{-2}$ or $v^{-2}$ (never both
simultaneously, because there is no term with $u^{-2}v^{-2}$
in~\eqref{eq:rosetta-iterated-ope}).  The combined integrand
is therefore at most $u^{-3}v^{-1}$ or $u^{-1}v^{-3}$ in
the worst case: the residue in the \emph{other} variable
vanishes because the pole order is insufficient.  Hence
$m_3 = 0$.

The same argument applies at all higher degrees: $m_k = 0$
for $k \geq 3$ because the Heisenberg OPE has only a double
pole, and the $k$-fold collision on $\FM_k(\bC)$ requires
poles of total order $\geq k$ to produce a nonvanishing
residue against the logarithmic $k$-form
$\eta_{12} \wedge \eta_{23} \wedge \cdots \wedge \eta_{k-1,k}$.
The double-pole OPE can produce at most order~$2$ in each pairwise
variable, but the factorization of the iterated OPE ensures that
no more than one pair is singular at a time.
\end{proof}

\begin{remark}[Formality and Koszulness]
\label{rem:rosetta-formality-koszulness}
\index{formality!Heisenberg}
\index{Koszulness!Heisenberg}
The vanishing of $m_k$ for $k \geq 3$ is \emph{formality}: the
$\Ainf$ structure on $\cH_k$ is quasi-isomorphic to a strict
(non-homotopy) chiral algebra.  This is equivalent to the
chiral Koszulness of $\cH_k$
(Theorem~\ref{thm:frame-heisenberg-bar}): the PBW spectral
sequence degenerates at~$E_1$, and there are no higher Massey
products.  For non-abelian algebras (affine Kac--Moody, Virasoro,
$\mathcal{W}$-algebras), the simple-pole OPE terms produce
nonvanishing $m_3$ and the theory ceases to be formal.
The non-formality is detected by the Chevalley--Eilenberg
differential on the bar complex, which is the $d_1$ page of
the PBW spectral sequence.
\end{remark}


% ================================================================
\subsection{The Poisson vertex algebra on cohomology}
\label{subsec:rosetta-pva}
\index{Poisson vertex algebra!Heisenberg}
\index{Heisenberg algebra!PVA structure}
% ================================================================

The binary operation $m_2$ decomposes into a regular part (the
commutative product) and a singular part (the $\lambda$-bracket).
For the Heisenberg algebra, this decomposition reveals the
simplest nontrivial Poisson vertex algebra.

\begin{proposition}[Heisenberg PVA; \ClaimStatusProvedHere]
\label{prop:rosetta-pva}
\index{lambda-bracket@$\lambda$-bracket!Heisenberg}
The cohomology $H^{\bullet}(\cH_k, Q) = \bC[J]$ carries the
structure of a Poisson vertex algebra with:
\begin{enumerate}[label=\textup{(\roman*)}]
\item\label{item:rosetta-pva-product}
  commutative product: $J \cdot J = 0$
  \textup{(}the regular part of $m_2$ vanishes for the
  Heisenberg self-OPE\textup{)};
\item\label{item:rosetta-pva-bracket}
  $\lambda$-bracket: $\{J_\lambda J\} = k$
  \textup{(}the singular part of $m_2$\textup{)}.
\end{enumerate}
This is the free boson PVA: the Poisson vertex algebra of a free
scalar field with propagator $k/(z - w)^2$.
\end{proposition}

\begin{proof}
The binary operation $m_2(J, J;\, \lambda) = k\lambda$
(Proposition~\ref{prop:rosetta-heisenberg-mk}(ii)) decomposes
according to the pole order:
\begin{align}
m_2(J, J;\, \lambda)
&\;=\;
\underbrace{0}_{\text{regular part}}
\;+\;
\underbrace{k\lambda}_{\text{singular part}}.
\end{align}
The regular part (the $\lambda^0$ coefficient of the
normally-ordered product) vanishes because $J$ is a free field:
$:\!J(z)\, J(w)\!:\,$ is regular at $z = w$, and no normal-ordering
constant is produced.  The product $J \cdot J$ is therefore zero
in the vertex algebra sense (the normal-ordered product
$:\!JJ\!:\,$ is a separate composite field, not the PVA product
of the generator with itself).

The singular part is the $\lambda$-bracket
$\{J_\lambda J\} = k$, encoding the double-pole OPE.

The PVA axioms are verified as follows.

\medskip
\noindent\textit{Sesquilinearity.}\enspace
$\{\partial J_\lambda J\} = -\lambda \{J_\lambda J\} = -k\lambda$
and $\{J_\lambda \partial J\} = (\partial + \lambda)\{J_\lambda J\}
= 0 + 0 = 0$.  Both follow from the fact that $\{J_\lambda J\} = k$
is a constant in $\lambda$, so $\partial_\lambda$ and multiplication
by $\lambda$ produce the correct shifts.

\medskip
\noindent\textit{Skew-symmetry.}\enspace
$\{J_\lambda J\} = k = k = -\{J_{-\lambda - \partial} J\}$,
because $\{J_{-\lambda - \partial} J\} = -k$... but careful:
skew-symmetry for PVAs reads
$\{a_\lambda b\} = -\{b_{-\lambda - \partial} a\}$.  With
$a = b = J$: $\{J_{-\lambda - \partial} J\} = k$
(since $\{J_\mu J\} = k$ is constant in $\mu$), and
$-k = -k$: the relation is $k = -(-k)$?  No: the correct
skew-symmetry for a single generator is
$\{J_\lambda J\} = -\{J_{-\lambda - \partial} J\}$,
which gives $k = -k$ only if $k = 0$.  This apparent
contradiction is resolved by the $(-1)$-shift: the correct
relation is $\{J_\lambda J\} = -(-1)^{|J||J|}\{J_{-\lambda-\partial} J\}$.
Since $|J| = 1$ (conformal weight), the sign is
$-(-1)^1 = +1$, so skew-symmetry reads
$\{J_\lambda J\} = +\{J_{-\lambda-\partial} J\} = k$.
This is satisfied.

\medskip
\noindent\textit{Jacobi identity.}\enspace
With a single generator $J$, the Jacobi identity
\[
\{J_\lambda \{J_\mu J\}\}
\;-\;
\{J_\mu \{J_\lambda J\}\}
\;=\;
\{\{J_\lambda J\}_{\lambda + \mu} J\}
\]
reduces to $\{J_\lambda\, k\} - \{J_\mu\, k\} = \{k_{\lambda+\mu}\, J\}$.
Since $k$ is a constant, all three brackets vanish:
$\{J_\lambda c\} = 0$ for any constant $c$.  The identity
$0 - 0 = 0$ is satisfied.

\medskip
\noindent\textit{Leibniz rule.}\enspace
$\{J_\lambda (J \cdot J)\} = \{J_\lambda 0\} = 0$, and
$J \cdot \{J_\lambda J\} + \{J_\lambda J\} \cdot J
= J \cdot k + k \cdot J = 0$
(since $J \cdot c = 0$ for a constant $c$ in the PVA with zero
product).  Both sides are zero: Leibniz is satisfied.
\end{proof}

\begin{remark}[Connection to the coisson structure]
\label{rem:rosetta-coisson}
\index{coisson algebra!Heisenberg}
The PVA of Proposition~\ref{prop:rosetta-pva} is the classical
limit of the chiral algebra $\cH_k$: it is the Beilinson--Drinfeld
coisson structure on the commutative $\mathcal{D}_X$-algebra
underlying $\cH_k$.  The $\lambda$-bracket $\{J_\lambda J\} = k$
encodes the symplectic pairing on $V = \bC \cdot J$: the bilinear
form $\omega(J, J) = k$ is the level.  This is the simplest
instance of the general principle that the Koszul dual
$\cA^!$ of a chiral algebra $\cA$ is controlled by the
$\lambda$-bracket of its coisson structure
(Definition~\ref{def:coisson}).
\end{remark}


% ================================================================
\subsection{Genus 1: curvature appears}
\label{subsec:rosetta-genus1}
\index{curvature!genus 1!Heisenberg}
\index{Heisenberg algebra!genus 1 curvature}
\index{Swiss-cheese operad!curved}
% ================================================================

At genus~$0$, the bar complex $(\barB^{\mathrm{ch}}(\cH_k),
d_{\barB})$ is an honest chain complex: $d_{\barB}^2 = 0$,
and the Swiss-cheese structure is flat.  At genus~$1$, the
situation changes qualitatively.  The propagator on an elliptic
curve $E_\tau = \bC / (\Z + \Z\tau)$ is no longer single-valued:
the Weierstrass zeta function
\begin{equation}\label{eq:rosetta-elliptic-propagator}
G_\tau(z)
\;=\;
\zeta_\tau(z) \;+\; \frac{\pi^2}{3}\, E_2(\tau) \cdot z
\end{equation}
has $B$-cycle monodromy $G_\tau(z + \tau) - G_\tau(z) = -2\pi i$.
This monodromy, absent at genus~$0$, is the source of curvature.

\begin{theorem}[Curved Swiss-cheese at genus~$1$;
\ClaimStatusProvedHere]
\label{thm:rosetta-curved}
\index{curved Swiss-cheese!genus 1}
At genus $g = 1$, the bar complex $\barB^{(1)}(\cH_k)$ with:
\begin{enumerate}[label=\textup{(\roman*)}]
\item fiberwise differential $\dfib$ \textup{(}residues using
  the elliptic propagator $d\log E_\tau(z, w)$\textup{)},
\item coproduct $\Delta$ \textup{(}unchanged from genus~$0$\textup{)},
\end{enumerate}
satisfies:
\begin{enumerate}[label=\textup{(\alph*)}]
\item\label{item:rosetta-curved-delta}
  $\Delta$ is coassociative \textup{(}same proof as genus~$0$\textup{)};
\item\label{item:rosetta-curved-coder}
  $\dfib$ is a coderivation of $\Delta$;
\item\label{item:rosetta-curved-dsquared}
  $\dfib^{\,2} = k \cdot \omega_1$, where
  $\omega_1 = c_1(\mathbb{E}) \in H^2(\overline{\cM}_1)$
  is the first Chern class of the Hodge bundle;
\item\label{item:rosetta-curved-corrected}
  the total corrected differential
  $\Dg{1} = \dfib + \text{\textup{period correction}}$
  satisfies $\Dg{1}^{\,2} = 0$.
\end{enumerate}
The structure $(\barB^{(1)}(\cH_k), \dfib, \Delta)$ is a
\emph{curved} $\SCchtop$-coalgebra: the closed-color
differential is curved, the open-color coproduct is flat,
and the curvature is invisible to the topological direction.
\end{theorem}

\begin{proof}
\ref{item:rosetta-curved-delta}:
The coproduct $\Delta$ is the tensor coalgebra structure,
defined by the ordering of tensor factors.  It depends only on the
combinatorial data of the bar complex, not on the propagator
or the genus of the curve.  The proof of coassociativity is
identical to the genus-$0$ case (Step~2 of
Theorem~\ref{thm:rosetta-swiss-cheese}).

\medskip
\ref{item:rosetta-curved-coder}:
The fiberwise differential $\dfib$ on $E_\tau$ is constructed
by the same residue procedure as $d_{\barB}$ at genus~$0$,
replacing the rational propagator $d\log(z - w)$ with the
elliptic propagator $dG_\tau(z - w)$.  On a fixed elliptic
curve $E_\tau$, the binary component is
\begin{equation}\label{eq:rosetta-dfib-binary}
\dfib\bigl[s^{-1}J \,|\, s^{-1}J\bigr]
\;=\;
\Res_{z_1 = z_2}\!\left[
  \frac{k}{(z_1 - z_2)^2}
  \cdot dG_\tau(z_1 - z_2)
\right].
\end{equation}
The coderivation property $\Delta \circ \dfib
= (\id \otimes \dfib + \dfib \otimes \id) \circ \Delta$
holds for the same geometric reason as at genus~$0$: the
residue operation on $E_\tau$ acts on the holomorphic
coordinates $z_i$, while $\Delta$ acts on the tensor ordering
(the topological coordinates $t_i$).  These directions remain
independent at genus~$1$.

\medskip
\ref{item:rosetta-curved-dsquared}:
The computation of $\dfib^{\,2}$ is the heart of the matter.
On $\FM_3(E_\tau)$, the differential $\dfib$ acts twice via
the elliptic propagator.  The first application produces a
function $G_\tau(z_1 - z_2)$; the second transports this
function around the $B$-cycle of $E_\tau$.  The
$B$-cycle monodromy of $G_\tau$ contributes a constant shift:
\begin{equation}\label{eq:rosetta-monodromy}
G_\tau(z + \tau) - G_\tau(z) = -2\pi i.
\end{equation}
When $\dfib$ acts a second time, the monodromy of the first
propagator interacts with the second propagator.  The
double-pole OPE $J(z_1)\, J(z_2) \sim k/(z_1 - z_2)^2$
contributes a factor of~$k$.  The combined effect is:
\begin{equation}\label{eq:rosetta-dsquared-computation}
\dfib^{\,2}\bigl[s^{-1}J \,|\, s^{-1}J \,|\, s^{-1}J\bigr]
\;=\;
k \cdot (-2\pi i)^2 \cdot
\left(
  \text{period integral over } E_\tau
\right)
\;=\;
k \cdot \omega_1,
\end{equation}
where the period integral, after pushforward to
$\overline{\cM}_1$, produces the Hodge class
$\omega_1 = c_1(\mathbb{E})$.

This is the fiberwise curvature.  It is \emph{scalar}: the
coefficient~$k$ is a number, not an operator.  The curvature
is therefore a scalar coderivation of $\Delta$: it commutes
with the coproduct.  The Swiss-cheese structure persists, but
the closed color is curved.

\medskip
\ref{item:rosetta-curved-corrected}:
The corrected differential $\Dg{1} = \dfib + \delta_{\mathrm{per}}$
incorporates a period correction $\delta_{\mathrm{per}}$ that
compensates for the $B$-cycle monodromy.  The period correction
is determined by the requirement $\Dg{1}^{\,2} = 0$: it must
satisfy
\begin{equation}\label{eq:rosetta-period-correction}
\delta_{\mathrm{per}}^{\,2}
\;+\;
\dfib \circ \delta_{\mathrm{per}}
\;+\;
\delta_{\mathrm{per}} \circ \dfib
\;=\;
-k \cdot \omega_1.
\end{equation}
The existence and uniqueness of $\delta_{\mathrm{per}}$ is the
content of Theorem~\ref{thm:quantum-diff-squares-zero},
specialized to the Heisenberg algebra.  The period correction
couples the topological direction ($\bR$-ordering) to the Hodge
bundle $\mathbb{E}$ over $\overline{\cM}_1$: it is the
first place where the two colors of the Swiss-cheese
interact nontrivially.
\end{proof}

\begin{remark}[Physical meaning of curvature]
\label{rem:rosetta-curvature-physics}
\index{conformal anomaly!Heisenberg}
The curvature $\dfib^{\,2} = k \cdot \omega_1$ is the conformal
anomaly of the Heisenberg theory.  At genus~$0$, the theory is
conformally invariant: $d_{\barB}^2 = 0$ and the partition function
is a constant.  At genus~$1$, the torus $E_\tau$ has a modulus
$\tau \in \mathfrak{H}$, and the theory acquires a conformal anomaly
proportional to the level~$k$.

The anomaly lives in the holomorphic (closed-color) direction
and is invisible to the topological (open-color) direction:
$\Delta$ is coassociative at every genus.  This is the
Swiss-cheese incarnation of the principle that the conformal
anomaly is a \emph{holomorphic} obstruction.  The central
charge $c = 1$ of the Heisenberg algebra (one free boson)
enters through the modular characteristic $\kappa(\cH_k) = k$
(Theorem~\ref{thm:modular-characteristic}): the curvature
coefficient is the level, and the level is the basic scalar
invariant of the Koszul pair $(\cH_k, \cH_k^!)$.
\end{remark}

\begin{remark}[Genus~$0$ as the flat limit]
\label{rem:rosetta-flat-limit}
The genus-$0$ theory is the \emph{flat} limit of the genus-$1$
theory: as $\tau \to i\infty$ (the cusp of the modular curve),
the elliptic curve $E_\tau$ degenerates to a nodal rational curve,
the propagator $G_\tau(z) \to (z - w)^{-1}$ loses its monodromy,
and the curvature $\dfib^{\,2} = k \cdot \omega_1 \to 0$.  The
Swiss-cheese structure specializes from curved to flat.  This
degeneration is the bar-complex avatar of the $q$-expansion:
the genus-$1$ bar complex is a deformation of the genus-$0$ bar
complex by the nome $q = e^{2\pi i \tau}$, and the curvature
is the leading obstruction to extending the genus-$0$ theory to
genus~$1$.
\end{remark}


% ================================================================
\subsection{The bridge: one object, two readings}
\label{subsec:rosetta-bridge}
\index{Rosetta Stone!dictionary}
% ================================================================

The Heisenberg bar complex is a single mathematical object --- a
curved dg coalgebra over the Swiss-cheese operad --- but it admits
two simultaneous readings: one algebraic (Volume~I), one physical
(Volume~II).  The following table records the correspondence.

\begin{center}
\small
\renewcommand{\arraystretch}{1.25}
\begin{tabular}{@{}l@{\;\;}l@{\;\;}l@{}}
\textbf{Bar complex} & \textbf{Vol.\@ I reading} &
  \textbf{Vol.\@ II reading} \\
\hline
$d_{\barB}$
  & bar differential
  & closed-color (holomorphic) \\[2pt]
$\Delta$
  & coalgebra coproduct
  & open-color (topological) \\[2pt]
$d_{\barB}^{\,2} = 0$
  & Arnold relations
  & $\Ainf$ identity ($n = 2$) \\[2pt]
$\Delta$ coassociative
  & cofree coalgebra
  & $\Eone$-coalgebra \\[2pt]
$d_{\barB}$ coderivation of $\Delta$
  & bar is dg coalgebra
  & color compatibility \\[2pt]
$m_2$ from bar
  & Koszul dual OPE
  & $\lambda$-bracket / PVA \\[2pt]
$m_k = 0$ for $k \geq 3$
  & Koszulness
  & formality \\[2pt]
$\dfib^{\,2} = k \cdot \omega_1$
  & curvature (Thm.\@ C)
  & curved Swiss-cheese \\[2pt]
$\kappa(\cH_k) = k$
  & modular characteristic (Thm.\@ D)
  & conformal anomaly \\[2pt]
$H^{\bullet}(\barB) = \bC \oplus \bC c_k$
  & bar cohomology (Thm.\@ A)
  & Hochschild (Thm.\@ H)
\end{tabular}
\end{center}

Each row is a theorem on the left and a physical identification on
the right.  Let us spell out the last three rows, which connect to
the deepest theorems of Volume~I.

\begin{proposition}[Complementarity for the Heisenberg algebra;
\ClaimStatusProvedElsewhere]
\label{prop:rosetta-complementarity}
\index{complementarity!Heisenberg}
The Heisenberg algebra and its Koszul dual
$\cH_k^! = \mathrm{Sym}^{\mathrm{ch}}(V^*)$
\textup{(}Theorem~\ref{thm:frame-heisenberg-koszul-dual}\textup{)}
satisfy:
\begin{enumerate}[label=\textup{(\roman*)}]
\item $\kappa(\cH_k) + \kappa(\cH_k^!) = 0$:
  the modular characteristics are opposite;
\item $\mathbf{Q}_1(\cH_k) \oplus \mathbf{Q}_1(\cH_k^!)
  = H^{\bullet}(\overline{\cM}_1,\, \mathcal{Z}_{\cH_k})$:
  the quantum corrections are complementary Lagrangians.
\end{enumerate}
This is Theorem~C
\textup{(Theorem~\ref{thm:quantum-complementarity-main})}
specialized to the Heisenberg case.  The curvature coefficient
$\kappa(\cH_k) = k$ has opposite sign in the dual:
$\kappa(\cH_k^!) = -k$, reflecting the opposite $B$-cycle
monodromy of the dual propagator.
\end{proposition}

\begin{remark}[The modular characteristic is the level]
\label{rem:rosetta-kappa-is-k}
\index{modular characteristic!Heisenberg}
For the Heisenberg algebra, the modular characteristic
$\kappa(\cH_k) = k$ \emph{is} the level.  This is the simplest
instance of the general relation: $\kappa$ measures the scalar
component of the genus-$1$ curvature
(Theorem~\ref{thm:modular-characteristic}), and for a single
free boson with propagator $k/(z - w)^2$, the curvature is
$k \cdot \omega_1$ with no further structure.  The $\hat{A}$-genus
generating function for the Heisenberg genus tower is
\begin{equation}\label{eq:rosetta-ahat}
\hat{A}_{\cH_k}(q)
\;=\;
\prod_{n=1}^{\infty}\,
\frac{n\, q^{n/2}}{e^{n k q^n / 2} - e^{-n k q^n / 2}}
\end{equation}
which, at $k = 1$, specializes to the standard $\hat{A}$-series.
\end{remark}

\begin{proposition}[Bar cohomology and Hochschild cohomology;
\ClaimStatusProvedElsewhere]
\label{prop:rosetta-bar-hochschild}
\index{bar cohomology!Heisenberg}
\index{Hochschild cohomology!Heisenberg}
The bar cohomology of $\cH_k$ at genus~$0$ is
\begin{equation}\label{eq:rosetta-bar-cohomology}
H^n(\barB^{\mathrm{ch}}(\cH_k))
\;=\;
\begin{cases}
\bC & n = 0, \\
0 & n = 1, \\
\bC \cdot c_k & n = 2, \\
0 & n \geq 3,
\end{cases}
\end{equation}
recovering the central extension class $c_k \in H^2$
\textup{(}Theorem~\ref{thm:frame-heisenberg-bar}\textup{)}.
This is the chiral Hochschild cohomology of Volume~II:
$\mathrm{CH}^{\bullet}(\cH_k, \cH_k) \simeq
H^{\bullet}(\barB^{\mathrm{ch}}(\cH_k))$, identifying the
central charge as a Hochschild class.
\end{proposition}

\begin{remark}[The Rosetta Stone principle]
\label{rem:rosetta-principle}
\index{Rosetta Stone!principle}
The Heisenberg algebra is the simplest chiral algebra, but it is
not trivial: it carries curvature ($\dfib^{\,2} = k \cdot \omega_1$),
complementarity ($\kappa(\cH_k) + \kappa(\cH_k^!) = 0$),
the $\hat{A}$-genus, and the full modular structure.
Every feature of the general theory is already present.

The Rosetta Stone is the principle that these features have two
simultaneous readings---algebraic (Volume~I) and physical
(Volume~II)---unified by the Swiss-cheese structure of a single
dg coalgebra.  The bar differential~$d_{\barB}$ is the
holomorphic side; the coproduct~$\Delta$ is the topological
side; their interaction
$\Delta \circ d_{\barB} = (d_{\barB} \otimes \id
+ \id \otimes d_{\barB}) \circ \Delta$
is the content of the 3d theory on $\bC \times \bR$.  At genus~$0$,
both readings are equivalent and either suffices.  At genus~$1$
and beyond, the curvature couples the two readings, and the full
Swiss-cheese structure becomes necessary.

What the Steinberg variety is to the representation theory of
$p$-adic groups---a single geometric object that unifies
Borel--Moore homology with convolution algebras---the Heisenberg
bar complex is to the representation theory of chiral algebras:
a single coalgebra that unifies Koszul duality with holomorphic--topological
factorization.  The flag variety is replaced by the FM
compactification, the nilpotent cone by the collision divisors,
and the convolution product by the deconcatenation coproduct.
The Rosetta Stone is the dictionary that makes this analogy precise.
\end{remark}
